\documentclass[tikz,border=14pt]{standalone}
\usepackage{xcolor}
\definecolor{bg}{HTML}{171D2B}
\definecolor{accent}{HTML}{F7B538}
\usetikzlibrary{angles,quotes}
\begin{document}
\pagecolor{bg}
\begin{tikzpicture}[scale=1.25, text=white, font=\small]
  \coordinate (A) at (2.6,3);
  \coordinate (B) at (0,0);
  \coordinate (C) at (6,0);
  \coordinate (P) at (-2.2,3);
  \coordinate (Q) at (8.2,3);
  % garis melalui A sejajar BC
  \draw[thick, accent] (P) -- (Q);
  % segitiga
  \draw[thick, white] (A) -- (B) -- (C) -- cycle;
  % sudut dalam segitiga
  \pic[draw, white, angle radius=0.6cm, angle eccentricity=1.7, "$\alpha$" {text=accent}] {angle = A--C--B};
  \pic[draw, white, angle radius=0.6cm, angle eccentricity=1.7, "$\beta$"  {text=accent}] {angle = C--B--A};
  % sudut di A: kiri = beta (sehadap dengan B), kanan = alpha (sehadap dengan C)
  \pic[draw, accent, angle radius=0.5cm, angle eccentricity=1.85, "$\beta$"  {text=accent}] {angle = P--A--B};
  \pic[draw, accent, angle radius=0.5cm, angle eccentricity=1.85, "$\alpha$" {text=accent}] {angle = C--A--Q};
  % titik dan nama
  \fill[accent] (A) circle (1.8pt) node[above=13pt, white] {$A$};
  \fill[accent] (B) circle (1.8pt) node[below left=-2pt, white] {$B$};
  \fill[accent] (C) circle (1.8pt) node[below right=-2pt, white] {$C$};
  \node[accent, above] at (7.0,3) {\scriptsize $\parallel BC$};
  \useasboundingbox (-2.8,-0.8) rectangle (8.8,4.0);
\end{tikzpicture}
\end{document}